\section{建康、南渡与一座受限的朝廷}
\label{sec:eastern-jin-establishment-southern-society}

永嘉元年，司马睿出镇建业。对于正在北方内战中失去支点的西晋宗室而言，这不是把一套完备朝廷搬到长江以南，而是先在一处既有的区域中心寻找可以接续诏令、官员与军队的关系网。王导等人参与联络江东士族、整顿东南军政，遂使司马睿的府署逐渐成为南方政治中心的雏形。后来以建康为名的这座城市，因而并非只是一座避难的城池：它同时是晋室、先到的官僚与当地精英必须共同使用、又各自试图影响的地点。

《晋书·元帝纪》与〈王导传〉保存了司马睿出镇和王导参与的基本线索，《建康实录》则提供都城视角下的地名与人物线索。三者都不足以把最初的结盟写成一次众望所归的迎接。谁在何时表示支持、府署究竟能够调动多少州郡，以及江东旧族如何看待北来的宗室，传世文字并没有同样清楚的答案。所谓“侨姓”与“江东旧族”也不是两个内部毫无差别的整体；把后来形成的门第格局倒投到三〇七年，只会掩盖建业政权最初的试探与协商。

\subsection{城破、渡江与可用资源的重新组合}

北方的灾难使这种试探变得急迫。三一一年洛阳陷落，怀帝被俘；三一三年怀帝被杀，长安另立的愍帝朝廷也在三一六年随着刘曜攻陷长安而终结。几次都城的失守并不自动规定每一个人的去向，却使原有的官僚任命、军队供给和家族居处失去共同的中心。南下者带来的不只是人口，也有官职资格、旧有从属关系、土地诉求和对北方局势各不相同的消息。建业的府署能够借此扩大，却也必须面对安置、征发和地方秩序的压力。

因此，“南渡”不应被想象成整齐的一次迁徙，更不能由正史中带有修辞色彩的死散、流亡数字推算出精确的人口总量。有人抵达江东，有人留在淮北、河南或关中，也有人在行进、停留和再迁徙之间改换依附。侨置州郡和侨居名目的出现，说明朝廷试图在籍贯记忆、任官资格与实际居住之间建立行政秩序；它们却不能单独证明某一地已经聚居了多少人，或当地社会已经接受了何种安排。南方朝廷的社会基础，正是在这种不均匀的流动、交涉与征发中形成，而不是由一个抽象的“北来士族”一次建立。

建兴元年祖逖率部北渡，在豫州、河南一带屯田并与地方力量结盟，显示建业一方并未放弃江北联系。可这类经营依赖将领本人的人际网络、能够获得的粮食与当地关系，并不等于一条由建业稳定供给、统一指挥的北伐战线。《晋书·祖逖传》与《资治通鉴》足以确认北渡、屯田和经营的轮廓；它们所载收复范围、兵粮数额和归附情形却不宜照录为精确地图。《世说新语》注所引《晋阳秋》可补充后世流传的线索，仍不能把人物声名替代为军政控制的证据。

\subsection{从晋王到皇帝}

长安陷落之后，三一七年司马睿在建康即晋王位，南方朝廷由区域府署转而公开承担晋室继承的名义。次年称帝、改元，东晋才以建康为都正式确立。帝号解决的是“谁可以代表晋”的迫切问题，并没有同时解决“谁能征发兵粮、约束军镇、处理流民”的问题。北方多政权并立，长江、淮河与江汉之间的道路又分别联系着不同的军镇和地方社会；建康所能直接触及的资源，远小于西晋全盛时洛阳所宣称的天下。

《晋书·元帝纪》、〈王导传〉与《资治通鉴》对即王、称帝的前后次序可以互相校核，《建康实录》也有助于辨认都城叙事的线索。不过，受禅礼仪究竟由哪些北方士人代表，南方地方社会以何种方式接受新皇统，均不能从这些以朝廷人物为中心的材料中直接得出。把三一八年写成旧秩序在江南无损地续存，或写成一个完全脱离北方的新国家，同样会失真。东晋首先是一个以晋朝名义重新组织的南方朝廷；它的财政、兵源和北方目标，始终受区域分裂所限制。

\subsection{平乱的军镇与王敦的压力}

这种限制很快显露在长江中游。大兴三年，荆湘地区的杜弢叛乱被王敦等人平定，朝廷重新取得该地区一部分军政控制。战争的结果不能简单化为“秩序恢复”：永嘉以后的人口流动、州郡征发与地方武装的竞争，都可能使不同群体卷入冲突。《晋书·杜弢传》与〈王敦传〉保留的是以朝廷和将领为中心的叙述，难以据此逐一判定流民、地方豪强及史书所称其他群体的参与程度。较为清楚的是，建康要维系长江中游，必须依靠能够自行调兵、筹给的军镇将领；平乱也因此增加了王敦的军事声望。

大兴四年，王敦以“清君侧”为名自武昌东下，迫使元帝调整中枢人事，并承认其军事实力。这不是一个只由个人好恶引起的宫廷插曲。东晋初建时，王氏的政治网络和荆州方面的兵力原是朝廷得以运转的条件；当元帝试图扶植近臣、重建较可控制的中枢时，外镇便把对任用和诏令的争议转化为武力压力。《晋书·元帝纪》、〈王敦传〉和〈王导传〉对各方的言辞、调停与责任分配并不完全一致，《资治通鉴》的连贯叙述也不能消除这些差异。可以把握的是，建康不得不让步，而皇帝对江上军镇的约束由此显著削弱。

永昌元年，王敦再度起兵，攻入建康，杀逐元帝近臣并控制朝政。城内战斗的经过、元帝每一步的应对和死者名单，在《晋书》诸传之间仍有需要核对之处；不必用戏剧化的宫门细节填补材料的空白。事件本身已经说明，建康虽是皇统所在地，却不能仅凭都城、礼仪和官号支配沿江的军事资源。王敦没有立即废帝，反而更能显示当时的结构：晋朝名义仍有价值，能够决定日常朝局的力量却可以来自都城之外。

\subsection{王敦死后，都城仍无定局}

永昌元年闰十一月，元帝去世，明帝司马绍继位。继承发生在王敦集团尚未解除威胁之时，遗诏如何安排、明帝实际能调动多少兵力、王氏内部是否同心，都不能由《晋书》诸传的事后叙述直接断定。可以把握的是，新皇帝必须依靠地方将领与王氏成员之间的分化来恢复建康的空间。太宁二年，王敦再举兵而病死途中，其部众在建康内外败散；这使朝廷暂时收回都城，却没有消除大族和军镇握有兵力、能够影响中枢的结构。

咸和二年，历阳将苏峻因拒绝征召而起兵，并与祖约威胁建康。庾亮等人试图削弱掌兵的流民将领，苏峻担忧失去兵权，这是冲突的可见条件；兵力数字、士族的态度与征召缘由的全部细节仍有争议。次年苏峻军沿江攻入建康，劫持成帝并控制宫城。城战、掠夺和幼帝的行动空间不宜写成能够逐项复原的场面，但陷都本身表明，王敦失败以后，朝廷仍没有建立一套能独立于外镇和流民军的防卫机制。各地军镇以勤王名义再度组织，既救援皇统，也重新分配了谁有资格在都城发言。

\subsection{北方的易手与南方的有限连续}

建康的危机同北方并行，并不受其直接指挥。三二八年，石勒在洛阳附近击败并俘获前赵皇帝刘曜；次年后赵攻取长安，前赵灭亡。战场位置、双方兵数和关中城内的处置各有异文，但河南与关中先后易手，使襄国成为北方最强的军事中心。对东晋而言，这不是收复旧都的机会自动到来，而是又一次提醒：北方的政权变化会重写可供流亡者、军镇和南方朝廷判断的道路，却不会自动补足建康的兵粮与命令。

\subsection{史料所照见的朝廷与未被照见的社会}

这一段年序以《晋书》的帝纪、列传和载记为主干。它成书于唐初，汇集较早史料，保存了司马睿、王导、祖逖、王敦、苏峻及北方政权行动的传世叙事；其人物褒贬、事后归责和宫廷秘闻却不能直接当作三、四世纪的现场记录。《资治通鉴》将这些事迹编入连续年代，特别适合核看都城陷落、即位与起兵的前后关联，但它是北宋的重新编排，不是独立的建康档案。《建康实录》保留都城和地方线索，也因成书较晚、传本与引据复杂，不能替代前两者。

城址、墓葬、墓志以及侨州郡的地方性记述，能够从另一种尺度提示城市空间、仕宦网络、籍贯自述和迁徙留下的痕迹；它们同样有发掘范围、墓主自我呈现、追赠官衔和传本层次的限制。一处墓葬不能代表整座建康，一条侨置记载也不能变成人口普查。本节因而只把它们作为校正正史视野的必要材料类型，而没有以考古或地方记录虚构三一七年拥立的现场、南渡人口的总数，或江东社会一致的态度。可确证的，是朝廷在建康重新竖起了晋室名号，并在迁徙、地方军镇和北方分裂的牵制下求取存续；其余那些无法由现存材料逐项回答的社会经验，应当保留为问题，而非用后来的整齐叙事加以填平。


% 年序补线：本节将既有账本中尚未初次落入正文的条目接入相应历史场景。
\subsection{两京失守：残余朝廷与南北道路}
记载中的光初二年（319年），前赵在\eventanchor{JH-0319-LIU-YAO-RENAMES-ZHAO}{刘曜在长安改汉国号为赵，以关中为中心重建政权身份}处留下刘曜在长安改汉国号为赵，以关中为中心重建政权身份的痕迹。 这里的资源与选择关系可从石勒在河北自立后，刘曜需要以关中行政与新国号凝聚原汉赵力量看出，结果使前赵与后赵的对峙使北方西东两大军事中心长期竞争。 材料的可说范围由《晋书·刘曜载记》；《资治通鉴·晋纪》；《十六国春秋》辑佚界定；改国号与长安中心地位可靠；改号时间、正统论述和前赵与后赵命名的后世区分有争议。

\eventanchor{JH-0319-SHI-LE-LATER-ZHAO}{石勒在襄国称赵王并建立独立政权，后赵由汉赵将领集团分出}所标出的石勒在襄国称赵王并建立独立政权，后赵由汉赵将领集团分出，发生在太兴二年（319年）的后赵。 它不是孤立转折：刘聪死后汉赵内部多变，石勒掌握河北军队和地方资源而不愿受平阳节制构成背景，北方形成汉赵、后赵及东晋并立格局，河北成为后赵长期基地则把压力送往下一阶段。 所据《晋书·石勒载记》；《资治通鉴·晋纪》；《十六国春秋》辑佚留下可核的线索，同时提醒读者称王、都襄国和与汉赵决裂可靠；部众构成、称号变化与控制范围有争议。

把\eventanchor{JH-LEDGER-032-1}{石勒在襄国称赵王并建立独立政权，后赵由汉赵将领集团分出}、\eventanchor{JH-LEDGER-032-2}{石勒在襄国称赵王并建立独立政权，后赵由汉赵将领集团分出}、\eventanchor{JH-LEDGER-032-3}{石勒在襄国称赵王并建立独立政权，后赵由汉赵将领集团分出}、\eventanchor{JH-LEDGER-032-4}{石勒在襄国称赵王并建立独立政权，后赵由汉赵将领集团分出}、\eventanchor{JH-LEDGER-032-5}{石勒在襄国称赵王并建立独立政权，后赵由汉赵将领集团分出}、\eventanchor{JH-LEDGER-032-6}{“石勒在襄国称赵王并建立独立政权，后赵由汉赵将领集团分出”的异源材料与影响范围的复核}放在一起读，才能看见太兴二年（319年）后赵所经历的石勒在襄国称赵王并建立独立政权，后赵由汉赵将领集团分出是一条事件链。 现有材料只能把这一变化放在连续年序中；前期条件、命令响应、地方执行、后续调整和异文问题是同一事件的不同读法，不能伪装成六场独立历史。 《晋书·石勒载记》；《资治通鉴·晋纪》；《十六国春秋》辑佚提供该簇的最低证据链；核心事件已有传世材料可供核对；各个观察面向的参与者、执行范围和时间先后仍须逐案核验。

光初二年（319年），前赵围绕\eventanchor{JH-LEDGER-033-1}{刘曜在长安改汉国号为赵，以关中为中心重建政权身份}、\eventanchor{JH-LEDGER-033-2}{刘曜在长安改汉国号为赵，以关中为中心重建政权身份}、\eventanchor{JH-LEDGER-033-3}{刘曜在长安改汉国号为赵，以关中为中心重建政权身份}、\eventanchor{JH-LEDGER-033-4}{刘曜在长安改汉国号为赵，以关中为中心重建政权身份}、\eventanchor{JH-LEDGER-033-5}{刘曜在长安改汉国号为赵，以关中为中心重建政权身份}、\eventanchor{JH-LEDGER-033-6}{“刘曜在长安改汉国号为赵，以关中为中心重建政权身份”的异源材料与影响范围的复核}留下的其实是同一事件簇：刘曜在长安改汉国号为赵，以关中为中心重建政权身份。 现有材料只能把这一变化放在连续年序中；前期条件、命令响应、地方执行、后续调整和异文问题是同一事件的不同读法，不能伪装成六场独立历史。 《晋书·刘曜载记》；《资治通鉴·晋纪》；《十六国春秋》辑佚提供该簇的最低证据链；核心事件已有传世材料可供核对；各个观察面向的参与者、执行范围和时间先后仍须逐案核验。

\subsection{建康初立：军镇、流民与北方竞争}
大兴三年（320年）的年序到\eventanchor{JH-0320-DU-TAO-DEFEATED}{杜弢领导的荆湘叛乱被王敦等镇压，东晋重新取得长江中游部分军政控制}时，东晋与荆湘地方集团须面对杜弢领导的荆湘叛乱被王敦等镇压，东晋重新取得长江中游部分军政控制。 行动者受永嘉后人口流动、州郡征发和地方军阀竞争使荆湘秩序恶化限制，随后王敦的军事声望上升，东晋对长江中游的控制更依赖强藩。 这里以《晋书·杜弢传》；《晋书·王敦传》；《资治通鉴·晋纪》为证据支点，且明确保留叛乱与其被平定可靠；流民、蛮獠与地方豪强的参与程度及伤亡有争议。

对于大兴三年（320年）的东晋与荆湘地方集团而言，\eventanchor{JH-LEDGER-034-1}{杜弢领导的荆湘叛乱被王敦等镇压，东晋重新取得长江中游部分军政控制}、\eventanchor{JH-LEDGER-034-2}{杜弢领导的荆湘叛乱被王敦等镇压，东晋重新取得长江中游部分军政控制}、\eventanchor{JH-LEDGER-034-3}{杜弢领导的荆湘叛乱被王敦等镇压，东晋重新取得长江中游部分军政控制}、\eventanchor{JH-LEDGER-034-4}{杜弢领导的荆湘叛乱被王敦等镇压，东晋重新取得长江中游部分军政控制}、\eventanchor{JH-LEDGER-034-5}{杜弢领导的荆湘叛乱被王敦等镇压，东晋重新取得长江中游部分军政控制}、\eventanchor{JH-LEDGER-034-6}{“杜弢领导的荆湘叛乱被王敦等镇压，东晋重新取得长江中游部分军政控制”的异源材料与影响范围的复核}都回到同一具体节点：杜弢领导的荆湘叛乱被王敦等镇压，东晋重新取得长江中游部分军政控制。 现有材料只能把这一变化放在连续年序中；前期条件、命令响应、地方执行、后续调整和异文问题是同一事件的不同读法，不能伪装成六场独立历史。 《晋书·杜弢传》；《晋书·王敦传》；《资治通鉴·晋纪》提供该簇的最低证据链；核心事件已有传世材料可供核对；各个观察面向的参与者、执行范围和时间先后仍须逐案核验。

大兴四年（321年），东晋的\eventanchor{JH-0321-WANG-DUN-FIRST-REBELLION}{王敦以清君侧为名自武昌东下，迫使元帝调整中枢人事并承认其军事优势}写下王敦以清君侧为名自武昌东下，迫使元帝调整中枢人事并承认其军事优势这一处具体变化。 它承接东晋初建依赖王氏军政网络，元帝试图培植刘隗、刁协等近臣触发强藩反弹，并使皇权对江上军镇的约束明显削弱，建康政治长期受外镇压力塑造。 这段叙述只越不过《晋书·元帝纪》；《晋书·王敦传》；《晋书·王导传》；《资治通鉴·晋纪》所能承载的范围，王敦起兵和朝廷让步可靠；王敦具体诉求、王导调停效果和各州响应有争议。

永昌元年（322年），\eventanchor{JH-0322-WANG-DUN-TAKES-JIANKANG}{王敦第二次起兵攻入建康，杀逐元帝近臣并控制朝政}并非抽象名称，而是东晋发生王敦第二次起兵攻入建康，杀逐元帝近臣并控制朝政的叙事入口。 从中可辨认出元帝试图重建禁军和限制王敦，引发荆州军政集团以武力进入都城的约束，及其对东晋皇权遭重创，王敦虽未废帝却成为决定朝局的强人的影响。 据《晋书·元帝纪》；《晋书·王敦传》；《晋书·周顗传》；《资治通鉴·晋纪》，可确认的是基本顺序和所列行动；王敦入建康及中枢清洗可靠；城内战斗、元帝应对和死者名单有争议。

在永昌元年闰十一月（323年），\eventanchor{JH-0323-YUAN-DEATH}{元帝去世，明帝司马绍即位，王敦集团与建康朝廷的对抗延续}把东晋置于元帝去世，明帝司马绍即位，王敦集团与建康朝廷的对抗延续的历史现场。 将其置回事件链，前提是王敦控制朝政后元帝政治空间有限，继承发生在强藩威胁未除之际，而后续处置面对明帝需要依靠地方将领与王氏分化来恢复都城控制。 材料来自《晋书·元帝纪》；《晋书·明帝纪》；《晋书·王敦传》；因此卒年与继位可靠；遗诏、明帝掌握军权的程度和王氏内部立场有争议。

记载中的太宁二年（324年），东晋在\eventanchor{JH-0324-WANG-DUN-COLLAPSE}{王敦再度举兵但病死于途中，其部众在建康内外被明帝集团击败}处留下王敦再度举兵但病死于途中，其部众在建康内外被明帝集团击败的痕迹。 这一节点将王敦试图以武力彻底取代建康中枢，但内部将领对僭越前景并不一致与东晋暂时收回都城控制，却未能消除大族和军镇对政治的结构性影响接合，却不预设两者之间只有一种因果。 所据《晋书·明帝纪》；《晋书·王敦传》；《晋书·温峤传》；《资治通鉴·晋纪》留下可核的线索，同时提醒读者王敦病死和叛军失败可靠；病因、部将倒戈次序与朝廷清算范围有争议。

记载中的太宁三年（325年），东晋与河南地方集团在\eventanchor{JH-0325-ZU-TI-DEATH}{祖逖去世，其北伐军和河南经营失去核心协调者}处留下祖逖去世，其北伐军和河南经营失去核心协调者的痕迹。 其代价落在祖逖的北方据点高度依赖个人威望与地方结盟，建康难以持续供给所限定的处境中，并以东晋对黄河以北的影响进一步萎缩，后赵得以扩张河南延续下去。 这里以《晋书·祖逖传》；《晋书·郗鉴传》；《资治通鉴·晋纪》为证据支点，且明确保留卒年与北伐经营受挫可靠；部众去向、继任安排和北方地方势力的反应有争议。

咸和二年（327年），\eventanchor{JH-0327-SU-JUN-REBELLION}{历阳将苏峻以拒绝征召为由起兵，联合祖约威胁建康}并非抽象名称，而是东晋发生历阳将苏峻以拒绝征召为由起兵，联合祖约威胁建康的叙事入口。 行动者受庾亮试图削弱掌兵的流民将领，苏峻担忧被解除兵权而先发制人限制，随后建康再度面临军事危机，幼帝与庾氏外戚的统治基础被动摇。 材料来自《晋书·成帝纪》；《晋书·苏峻传》；《晋书·庾亮传》；《资治通鉴·晋纪》；因此苏峻、祖约起兵可靠；征召缘由、兵力数字和士族对叛军的态度有争议。

太和三年（328年）的年序到\eventanchor{JH-0328-SHI-LE-CAPTURES-LIU-YAO}{石勒在洛阳附近击败并俘获前赵皇帝刘曜，前赵主力由此崩解}时，后赵与前赵须面对石勒在洛阳附近击败并俘获前赵皇帝刘曜，前赵主力由此崩解。 这一处理回应前赵与后赵长期争夺河南，刘曜试图以洛阳为支点压制石勒，同时改变了后赵夺得河南与前赵皇帝，北方权力天平迅速向襄国倾斜。 《晋书·石勒载记》；《晋书·刘曜载记》；《资治通鉴·晋纪》构成最低证据链，故刘曜被俘和战役结局可靠；战场位置、双方兵数及刘曜战败原因有争议。

在咸和三年（328年）这一个节点，东晋通过\eventanchor{JH-0328-SU-JUN-TAKES-JIANKANG}{苏峻军攻入建康，劫持成帝并控制宫城和朝廷}显出苏峻军攻入建康，劫持成帝并控制宫城和朝廷的压力。 此处不把时间相邻当作因果；能追到的是朝廷军队缺乏协调，苏峻凭历阳水军沿江迅速推进，以及其后东晋中枢短暂瘫痪，各地军镇以勤王名义重新组织政治联盟。 《晋书·成帝纪》；《晋书·苏峻传》；《晋书·温峤传》；《建康实录》保留了这条年序，然而陷都、劫帝和宫城控制可靠；城战细节、掠夺程度与成帝行动空间有争议。

太和四年（329年），后赵与前赵的\eventanchor{JH-0329-LATER-ZHAO-DESTROYS-FORMER-ZHAO}{石勒攻取长安并处置刘曜之子刘熙，前赵政权灭亡}使石勒攻取长安并处置刘曜之子刘熙，前赵政权灭亡成为可追索的行动。 这一处理回应刘曜被俘后前赵指挥系统瓦解，关中难以获得有效援军，同时改变了后赵势力进入关中，石勒成为北方最强的军事统治者。 据《晋书·石勒载记》；《晋书·刘曜载记》；《资治通鉴·晋纪》，可确认的是基本顺序和所列行动；前赵灭亡和长安失守可靠；城内抵抗、处置过程和关中地方反应有争议。

\subsection{江上军镇：都城防卫的代价}
建平元年（330年），后赵的\eventanchor{JH-0330-SHI-LE-EMPEROR}{石勒在襄国称帝，后赵以河北为核心建立更完整的官僚与军政体系}写下石勒在襄国称帝，后赵以河北为核心建立更完整的官僚与军政体系这一处具体变化。 这里的资源与选择关系可从消灭前赵后石勒拥有河北、河南和关中部分资源，需要以帝号整合文武集团看出，结果使后赵成为北方主要政权，其征发与城镇建设深刻影响黄河流域社会。 材料来自《晋书·石勒载记》；《十六国春秋》辑佚；《资治通鉴·晋纪》；因此称帝与建平改元可靠；官制仿效程度、地方执行和羯胡—汉人官僚关系有争议。

建平四年（333年），后赵的\eventanchor{JH-0333-SHI-LE-DEATH}{石勒去世，太子石弘继位，石虎掌握军政实权}使石勒去世，太子石弘继位，石虎掌握军政实权成为可追索的行动。 由此能看到的是石勒建立的统治依赖个人威望与军事集团，继承人缺乏独立兵权与石虎的夺权使后赵进入高压扩张与内部猜忌并行的阶段之间的条件关系，而非事后必然论。 《晋书·石勒载记》；《晋书·石虎载记》；《资治通鉴·晋纪》构成最低证据链，故石勒卒、石弘继位与石虎专权可靠；遗诏内容、石弘自主空间和宗室军权分配有争议。

玉衡二十四年（334年），成汉的\eventanchor{JH-0334-LI-XIONG-DEATH}{李雄去世，李班继位后旋即被李期夺权，成汉内部继承冲突公开化}使李雄去世，李班继位后旋即被李期夺权，成汉内部继承冲突公开化成为可追索的行动。 它承接李雄长期统治下未能建立无争议的继承秩序，宗室各支均掌有军政资源，并使成汉从相对稳定转入频繁宫廷冲突，削弱其抵御东晋北伐的能力。 年序和核心行动以《晋书·李雄载记》；《晋书·李期载记》；《华阳国志·李特雄期寿势志》为据；李雄卒、李班继位与李期夺权可靠；遗命、宗室派系和成都官僚的作用有争议。

记载中的建武元年（334年），后赵在\eventanchor{JH-0334-SHI-HU-COUP}{石虎废杀石弘，自立为天王，后赵皇位转入其支系}处留下石虎废杀石弘，自立为天王，后赵皇位转入其支系的痕迹。 从中可辨认出石虎已控制禁军与外军，石弘无法制衡其叔父的军事网络的约束，及其对后赵政治更趋个人化和暴力化，继承危机被暂时压制而非解决的影响。 这项判断依托《晋书·石虎载记》；《晋书·石勒载记下》；《资治通鉴·晋纪》，并受限于废立与石虎掌权可靠；石弘死亡方式、宗室支持和称号礼仪有争议。

在咸康元年（335年），\eventanchor{JH-0335-MURONG-HUANG-YAN}{慕容皝在辽西称燕王，前燕政权以棘城为中心扩展}把前燕置于慕容皝在辽西称燕王，前燕政权以棘城为中心扩展的历史现场。 行动者受慕容部在辽西积累军事实力，后赵的北方统治也难以有效深入东北限制，随后前燕成为东北重要政权，参与后赵、东晋与高句丽之间的区域竞争。 材料的可说范围由《晋书·慕容皝载记》；《资治通鉴·晋纪》；《十六国春秋》辑佚界定；称燕王与前燕政权形成大体可靠；受晋册封关系、部族联盟和疆域范围有争议。

在咸康二年（336年）这一个节点，前燕与段部通过\eventanchor{JH-0336-MURONG-HUANG-DEFEATS-DUAN}{慕容皝击败段辽并整合辽西部众，前燕的军事与人口资源扩大}显出慕容皝击败段辽并整合辽西部众，前燕的军事与人口资源扩大的压力。 从中可辨认出辽西诸部长期竞争牧地、城郭和对后赵的关系，慕容皝寻求消除侧翼威胁的约束，及其对前燕强化对辽西通道的掌控，为后续与后赵和高句丽作战提供兵源的影响。 以《晋书·慕容皝载记》；《资治通鉴·晋纪》；《十六国春秋》辑佚交叉检验，能够确认的止于此；战胜段部和吸纳部众可靠；段部降附规模、部族身份和迁徙路线有争议。

咸康四年（338年），\eventanchor{JH-0338-LI-SHOU-USURPS-CHENGHAN}{李寿废李期自立并改国号为汉，成汉政权的宗室结构再次重组}并非抽象名称，而是成汉发生李寿废李期自立并改国号为汉，成汉政权的宗室结构再次重组的叙事入口。 它承接李雄死后继承冲突未止，李寿凭军权和宗室身份夺取成都，并使政权虽暂获统一但统治正当性受损，为东晋西征创造机会。 《晋书·李寿载记》；《华阳国志·李特雄期寿势志》；《资治通鉴·晋纪》可互校时间位置与主体，但废立、李寿称帝和国号变化可靠；李期罪状、成都官僚态度和改号的实际使用有争议。

沿着咸康四年（338年），\eventanchor{JH-0338-SHI-HU-ATTACKS-YAN}{石虎大军进攻前燕，慕容皝依托棘城及周边地形击退来军}让后赵与前燕的石虎大军进攻前燕，慕容皝依托棘城及周边地形击退来军进入连续叙事。 由此能看到的是后赵试图遏制慕容部扩张，前燕则需保住辽西的政治中心与后赵未能消灭前燕，东北独立政权得以继续发展之间的条件关系，而非事后必然论。 这段叙述只越不过《晋书·石虎载记》；《晋书·慕容皝载记》；《资治通鉴·晋纪》所能承载的范围，后赵进攻与前燕守住核心地区可靠；兵额、战场细节和战损数字有争议。



\subsection{335—344年：材料所见的连续压力}
建康须同时安置北来人与旧有地方关系，北方诸国的更替也影响水路与消息；材料只留下密度不一的户籍、施主与军镇痕迹。

在公元335年（年序桥接）的并行行动者之间，\eventanchor{JH-DENS-335-A}{长江中游军府、北方诸政权与益州之间的战争和继承竞争，使道路、屯田和地方征发持续成为节点}将长江中游军府、北方诸政权与益州之间的战争和继承竞争，使道路、屯田和地方征发持续成为节点落到时间位置，\eventanchor{JH-DENS-335-B}{东晋、后赵与成汉的关隘、渡口、军镇和水陆道路的可通行性继续约束使节、军队和物资的行动}从东晋、后赵与成汉的关隘、渡口、军镇和水陆道路的可通行性继续约束使节、军队和物资的行动补足资源与空间的视角。 如此处理不是用空白填满史实，而是避免后来的转折抹去本年仍在运作的条件。 《晋书》成康穆诸帝纪；《晋书》载记；《资治通鉴·晋纪》使前一线索可追，本纪与编年材料主要确证政权年序和精英行动；对基层执行、疆界和人口数量的沉默不得以叙事补足；《华阳国志》；墓志、城址和军镇考古限制后一线索的说法，纪传中的行军路线、兵数与控制范围常被简化或夸饰，地理材料只能提供有限交叉。

公元336年（年序桥接），\eventanchor{JH-DENS-336-A}{长江中游军府、北方诸政权与益州之间的战争和继承竞争，使道路、屯田和地方征发持续成为节点}和\eventanchor{JH-DENS-336-B}{东晋、后赵与成汉的朝廷和州郡围绕户籍、田租、军粮与转运的安排进入材料}共同避免把东晋、后赵与成汉压缩为单一都城事件：前者关注长江中游军府、北方诸政权与益州之间的战争和继承竞争，使道路、屯田和地方征发持续成为节点，后者关注东晋、后赵与成汉的朝廷和州郡围绕户籍、田租、军粮与转运的安排进入材料。 年序因此不把大事件当作一切的终点，而让制度、交通和社会关系继续承受压力。 《晋书》成康穆诸帝纪；《晋书》载记；《资治通鉴·晋纪》使前一线索可追，本纪与编年材料主要确证政权年序和精英行动；对基层执行、疆界和人口数量的沉默不得以叙事补足；《华阳国志》；墓志、城址和军镇考古限制后一线索的说法，志书、诏令与本纪可显示制度设计或战争需求，不能直接换算赋额、仓储或实际负担。

把公元337年（年序桥接）放回连续叙事，\eventanchor{JH-DENS-337-A}{长江中游军府、北方诸政权与益州之间的战争和继承竞争，使道路、屯田和地方征发持续成为节点}保留长江中游军府、北方诸政权与益州之间的战争和继承竞争，使道路、屯田和地方征发持续成为节点，\eventanchor{JH-DENS-337-B}{东晋、后赵与成汉的流民、侨置户、军镇居民或地方家族的安置与编户问题构成社会背景}让东晋、后赵与成汉的流民、侨置户、军镇居民或地方家族的安置与编户问题构成社会背景继续约束东晋、后赵与成汉的行动。 如此处理不是用空白填满史实，而是避免后来的转折抹去本年仍在运作的条件。 两种材料不能互相替代：《晋书》成康穆诸帝纪；《晋书》载记；《资治通鉴·晋纪》所能确认的有限，本纪与编年材料主要确证政权年序和精英行动；对基层执行、疆界和人口数量的沉默不得以叙事补足；《华阳国志》；墓志、城址和军镇考古同样只照见一部分，墓志、家族文本和侨置名目各自有选择性，不能据此统计迁徙规模或概括整体社会态度。

\eventanchor{JH-DENS-338-A}{长江中游军府、北方诸政权与益州之间的战争和继承竞争，使道路、屯田和地方征发持续成为节点}在公元338年（年序桥接）提示长江中游军府、北方诸政权与益州之间的战争和继承竞争，使道路、屯田和地方征发持续成为节点仍未结案；\eventanchor{JH-DENS-338-B}{东晋、后赵与成汉的寺院、僧侣、道士和施主网络在材料中连接都城、州郡与边地}以东晋、后赵与成汉的寺院、僧侣、道士和施主网络在材料中连接都城、州郡与边地说明东晋、后赵与成汉的选择还受具体条件限制。 这使读者能把下一年的军事、行政或社会变化放回尚未消散的限制中。 就可说范围而言，前一判断止于《晋书》成康穆诸帝纪；《晋书》载记；《资治通鉴·晋纪》，本纪与编年材料主要确证政权年序和精英行动；对基层执行、疆界和人口数量的沉默不得以叙事补足；后一判断止于《华阳国志》；墓志、城址和军镇考古，僧传、道教文本、造像记与遗址的成书或营造年代须分开，个案不能量化信仰或政策落实。

公元339年（年序桥接），东晋、后赵与成汉的\eventanchor{JH-DENS-339-A}{长江中游军府、北方诸政权与益州之间的战争和继承竞争，使道路、屯田和地方征发持续成为节点}把长江中游军府、北方诸政权与益州之间的战争和继承竞争，使道路、屯田和地方征发持续成为节点留在年序中；\eventanchor{JH-DENS-339-B}{东晋、后赵与成汉的关隘、渡口、军镇和水陆道路的可通行性继续约束使节、军队和物资的行动}则从东晋、后赵与成汉的关隘、渡口、军镇和水陆道路的可通行性继续约束使节、军队和物资的行动观察同一年的约束。 这里保留的是可核的连续性，而非对基层反应、疆界或人数的无据补写。 两种材料不能互相替代：《晋书》成康穆诸帝纪；《晋书》载记；《资治通鉴·晋纪》所能确认的有限，本纪与编年材料主要确证政权年序和精英行动；对基层执行、疆界和人口数量的沉默不得以叙事补足；《华阳国志》；墓志、城址和军镇考古同样只照见一部分，纪传中的行军路线、兵数与控制范围常被简化或夸饰，地理材料只能提供有限交叉。

公元340年（年序桥接），东晋、后赵与成汉的\eventanchor{JH-DENS-340-A}{长江中游军府、北方诸政权与益州之间的战争和继承竞争，使道路、屯田和地方征发持续成为节点}把长江中游军府、北方诸政权与益州之间的战争和继承竞争，使道路、屯田和地方征发持续成为节点留在年序中；\eventanchor{JH-DENS-340-B}{东晋、后赵与成汉的朝廷和州郡围绕户籍、田租、军粮与转运的安排进入材料}则从东晋、后赵与成汉的朝廷和州郡围绕户籍、田租、军粮与转运的安排进入材料观察同一年的约束。 年序因此不把大事件当作一切的终点，而让制度、交通和社会关系继续承受压力。 《晋书》成康穆诸帝纪；《晋书》载记；《资治通鉴·晋纪》只能确认前一观察的基本年序，本纪与编年材料主要确证政权年序和精英行动；对基层执行、疆界和人口数量的沉默不得以叙事补足；《华阳国志》；墓志、城址和军镇考古为后一观察提供交叉线索，志书、诏令与本纪可显示制度设计或战争需求，不能直接换算赋额、仓储或实际负担。

\eventanchor{JH-DENS-341-A}{长江中游军府、北方诸政权与益州之间的战争和继承竞争，使道路、屯田和地方征发持续成为节点}在公元341年（年序桥接）提示长江中游军府、北方诸政权与益州之间的战争和继承竞争，使道路、屯田和地方征发持续成为节点仍未结案；\eventanchor{JH-DENS-341-B}{东晋、后赵与成汉的流民、侨置户、军镇居民或地方家族的安置与编户问题构成社会背景}以东晋、后赵与成汉的流民、侨置户、军镇居民或地方家族的安置与编户问题构成社会背景说明东晋、后赵与成汉的选择还受具体条件限制。 这两项观察不宣称另有一场未载战争或政策，只把已有压力的时间位置和可见面向分开。 前者依据《晋书》成康穆诸帝纪；《晋书》载记；《资治通鉴·晋纪》；本纪与编年材料主要确证政权年序和精英行动；对基层执行、疆界和人口数量的沉默不得以叙事补足。后者参照《华阳国志》；墓志、城址和军镇考古；墓志、家族文本和侨置名目各自有选择性，不能据此统计迁徙规模或概括整体社会态度。

\eventanchor{JH-DENS-342-A}{长江中游军府、北方诸政权与益州之间的战争和继承竞争，使道路、屯田和地方征发持续成为节点}和\eventanchor{JH-DENS-342-B}{东晋、后赵与成汉的寺院、僧侣、道士和施主网络在材料中连接都城、州郡与边地}使公元342年（年序桥接）的东晋、后赵与成汉既保留长江中游军府、北方诸政权与益州之间的战争和继承竞争，使道路、屯田和地方征发持续成为节点，也不略去东晋、后赵与成汉的寺院、僧侣、道士和施主网络在材料中连接都城、州郡与边地。 它们将政治时间同资源和地方网络并置，却不把并置误作已被证实的直接因果。 就可说范围而言，前一判断止于《晋书》成康穆诸帝纪；《晋书》载记；《资治通鉴·晋纪》，本纪与编年材料主要确证政权年序和精英行动；对基层执行、疆界和人口数量的沉默不得以叙事补足；后一判断止于《华阳国志》；墓志、城址和军镇考古，僧传、道教文本、造像记与遗址的成书或营造年代须分开，个案不能量化信仰或政策落实。

在公元343年（年序桥接），\eventanchor{JH-DENS-343-A}{长江中游军府、北方诸政权与益州之间的战争和继承竞争，使道路、屯田和地方征发持续成为节点}所见的长江中游军府、北方诸政权与益州之间的战争和继承竞争，使道路、屯田和地方征发持续成为节点与\eventanchor{JH-DENS-343-B}{东晋、后赵与成汉的关隘、渡口、军镇和水陆道路的可通行性继续约束使节、军队和物资的行动}所见的东晋、后赵与成汉的关隘、渡口、军镇和水陆道路的可通行性继续约束使节、军队和物资的行动并行发生在东晋、后赵与成汉的历史空间。 年序因此不把大事件当作一切的终点，而让制度、交通和社会关系继续承受压力。 材料边界须分别保留：《晋书》成康穆诸帝纪；《晋书》载记；《资治通鉴·晋纪》与《华阳国志》；墓志、城址和军镇考古所见不同，且本纪与编年材料主要确证政权年序和精英行动；对基层执行、疆界和人口数量的沉默不得以叙事补足；纪传中的行军路线、兵数与控制范围常被简化或夸饰，地理材料只能提供有限交叉。

在公元344年（年序桥接）的并行行动者之间，\eventanchor{JH-DENS-344-A}{长江中游军府、北方诸政权与益州之间的战争和继承竞争，使道路、屯田和地方征发持续成为节点}将长江中游军府、北方诸政权与益州之间的战争和继承竞争，使道路、屯田和地方征发持续成为节点落到时间位置，\eventanchor{JH-DENS-344-B}{东晋、后赵与成汉的朝廷和州郡围绕户籍、田租、军粮与转运的安排进入材料}从东晋、后赵与成汉的朝廷和州郡围绕户籍、田租、军粮与转运的安排进入材料补足资源与空间的视角。 两条线索共同说明，政权变化之外，地方社会和空间条件仍在塑造行动者可以选择的办法。 《晋书》成康穆诸帝纪；《晋书》载记；《资治通鉴·晋纪》使前一线索可追，本纪与编年材料主要确证政权年序和精英行动；对基层执行、疆界和人口数量的沉默不得以叙事补足；《华阳国志》；墓志、城址和军镇考古限制后一线索的说法，志书、诏令与本纪可显示制度设计或战争需求，不能直接换算赋额、仓储或实际负担。


《晋书》、载记、《资治通鉴》和建康墓志可校年代；单墓不能代表江南社会。
从公元341年（材料补点）的记录看，东晋江南经由\eventanchor{JH-DENS-341-X}{墓志与家族年序补点：王兴之夫妇墓志的纪年为建康地区家族、官职与葬制研究提供一个可核的时间节点，不可据单墓外推江南社会。}处理墓志与家族年序补点：王兴之夫妇墓志的纪年为建康地区家族、官职与葬制研究提供一个可核的时间节点，不可据单墓外推江南社会。 由此能看到的是墓志把个人家族的自我呈现固定在年号与地点之中。与为与建康朝廷年序相对照的地方材料提供锚点。之间的条件关系，而非事后必然论。 材料的可说范围由《王兴之夫妇墓志》；《晋书·成帝纪》；建康地区墓葬研究界定；墓志可确证其铭文、纪年和自我表述；官职、亲属关系或地域代表性仍须与其他材料核对。



\subsection{345—354年：材料所见的连续压力}
石赵后期与江南朝廷并非两条平行线：北方重组改变人群、道路和军粮的去向，南方也要选择资源投向州郡、侨置和防线。

沿着公元345年（年序桥接），东晋、后赵与成汉的\eventanchor{JH-DENS-345-A}{长江中游军府、北方诸政权与益州之间的战争和继承竞争，使道路、屯田和地方征发持续成为节点}指向长江中游军府、北方诸政权与益州之间的战争和继承竞争，使道路、屯田和地方征发持续成为节点，\eventanchor{JH-DENS-345-B}{东晋、后赵与成汉的流民、侨置户、军镇居民或地方家族的安置与编户问题构成社会背景}则把东晋、后赵与成汉的流民、侨置户、军镇居民或地方家族的安置与编户问题构成社会背景接入同一条年序。 它们将政治时间同资源和地方网络并置，却不把并置误作已被证实的直接因果。 前者依据《晋书》成康穆诸帝纪；《晋书》载记；《资治通鉴·晋纪》；本纪与编年材料主要确证政权年序和精英行动；对基层执行、疆界和人口数量的沉默不得以叙事补足。后者参照《华阳国志》；墓志、城址和军镇考古；墓志、家族文本和侨置名目各自有选择性，不能据此统计迁徙规模或概括整体社会态度。

把公元346年（年序桥接）放回连续叙事，\eventanchor{JH-DENS-346-A}{长江中游军府、北方诸政权与益州之间的战争和继承竞争，使道路、屯田和地方征发持续成为节点}保留长江中游军府、北方诸政权与益州之间的战争和继承竞争，使道路、屯田和地方征发持续成为节点，\eventanchor{JH-DENS-346-B}{东晋、后赵与成汉的寺院、僧侣、道士和施主网络在材料中连接都城、州郡与边地}让东晋、后赵与成汉的寺院、僧侣、道士和施主网络在材料中连接都城、州郡与边地继续约束东晋、后赵与成汉的行动。 这里保留的是可核的连续性，而非对基层反应、疆界或人数的无据补写。 前者依据《晋书》成康穆诸帝纪；《晋书》载记；《资治通鉴·晋纪》；本纪与编年材料主要确证政权年序和精英行动；对基层执行、疆界和人口数量的沉默不得以叙事补足。后者参照《华阳国志》；墓志、城址和军镇考古；僧传、道教文本、造像记与遗址的成书或营造年代须分开，个案不能量化信仰或政策落实。

公元347年（年序桥接），\eventanchor{JH-DENS-347-A}{长江中游军府、北方诸政权与益州之间的战争和继承竞争，使道路、屯田和地方征发持续成为节点}和\eventanchor{JH-DENS-347-B}{东晋、后赵与成汉的关隘、渡口、军镇和水陆道路的可通行性继续约束使节、军队和物资的行动}共同避免把东晋、后赵与成汉压缩为单一都城事件：前者关注长江中游军府、北方诸政权与益州之间的战争和继承竞争，使道路、屯田和地方征发持续成为节点，后者关注东晋、后赵与成汉的关隘、渡口、军镇和水陆道路的可通行性继续约束使节、军队和物资的行动。 它们将政治时间同资源和地方网络并置，却不把并置误作已被证实的直接因果。 《晋书》成康穆诸帝纪；《晋书》载记；《资治通鉴·晋纪》使前一线索可追，本纪与编年材料主要确证政权年序和精英行动；对基层执行、疆界和人口数量的沉默不得以叙事补足；《华阳国志》；墓志、城址和军镇考古限制后一线索的说法，纪传中的行军路线、兵数与控制范围常被简化或夸饰，地理材料只能提供有限交叉。

公元348年（年序桥接），东晋、前燕与前秦的\eventanchor{JH-DENS-348-A}{前燕、前秦与东晋围绕河北、关中和淮汉走廊轮番调兵，扩张依赖州郡供给与地方合作}把前燕、前秦与东晋围绕河北、关中和淮汉走廊轮番调兵，扩张依赖州郡供给与地方合作留在年序中；\eventanchor{JH-DENS-348-B}{东晋、前燕与前秦的朝廷和州郡围绕户籍、田租、军粮与转运的安排进入材料}则从东晋、前燕与前秦的朝廷和州郡围绕户籍、田租、军粮与转运的安排进入材料观察同一年的约束。 年序因此不把大事件当作一切的终点，而让制度、交通和社会关系继续承受压力。 证据不许把这两个观察写成全域事实：《晋书》哀海西简文诸帝纪；《晋书》载记；《资治通鉴·晋纪》受本纪与编年材料主要确证政权年序和精英行动；对基层执行、疆界和人口数量的沉默不得以叙事补足。限制，《水经注》；墓志、城址与十六国史辑佚亦受志书、诏令与本纪可显示制度设计或战争需求，不能直接换算赋额、仓储或实际负担限制。

\eventanchor{JH-DENS-349-A}{前燕、前秦与东晋围绕河北、关中和淮汉走廊轮番调兵，扩张依赖州郡供给与地方合作}在公元349年（年序桥接）提示前燕、前秦与东晋围绕河北、关中和淮汉走廊轮番调兵，扩张依赖州郡供给与地方合作仍未结案；\eventanchor{JH-DENS-349-B}{东晋、前燕与前秦的流民、侨置户、军镇居民或地方家族的安置与编户问题构成社会背景}以东晋、前燕与前秦的流民、侨置户、军镇居民或地方家族的安置与编户问题构成社会背景说明东晋、前燕与前秦的选择还受具体条件限制。 如此处理不是用空白填满史实，而是避免后来的转折抹去本年仍在运作的条件。 材料边界须分别保留：《晋书》哀海西简文诸帝纪；《晋书》载记；《资治通鉴·晋纪》与《水经注》；墓志、城址与十六国史辑佚所见不同，且本纪与编年材料主要确证政权年序和精英行动；对基层执行、疆界和人口数量的沉默不得以叙事补足；墓志、家族文本和侨置名目各自有选择性，不能据此统计迁徙规模或概括整体社会态度。

公元350年（年序桥接），\eventanchor{JH-DENS-350-A}{前燕、前秦与东晋围绕河北、关中和淮汉走廊轮番调兵，扩张依赖州郡供给与地方合作}和\eventanchor{JH-DENS-350-B}{东晋、前燕与前秦的寺院、僧侣、道士和施主网络在材料中连接都城、州郡与边地}共同避免把东晋、前燕与前秦压缩为单一都城事件：前者关注前燕、前秦与东晋围绕河北、关中和淮汉走廊轮番调兵，扩张依赖州郡供给与地方合作，后者关注东晋、前燕与前秦的寺院、僧侣、道士和施主网络在材料中连接都城、州郡与边地。 这使读者能把下一年的军事、行政或社会变化放回尚未消散的限制中。 证据不许把这两个观察写成全域事实：《晋书》哀海西简文诸帝纪；《晋书》载记；《资治通鉴·晋纪》受本纪与编年材料主要确证政权年序和精英行动；对基层执行、疆界和人口数量的沉默不得以叙事补足。限制，《水经注》；墓志、城址与十六国史辑佚亦受僧传、道教文本、造像记与遗址的成书或营造年代须分开，个案不能量化信仰或政策落实限制。

把公元351年（年序桥接）放回连续叙事，\eventanchor{JH-DENS-351-A}{前燕、前秦与东晋围绕河北、关中和淮汉走廊轮番调兵，扩张依赖州郡供给与地方合作}保留前燕、前秦与东晋围绕河北、关中和淮汉走廊轮番调兵，扩张依赖州郡供给与地方合作，\eventanchor{JH-DENS-351-B}{东晋、前燕与前秦的关隘、渡口、军镇和水陆道路的可通行性继续约束使节、军队和物资的行动}让东晋、前燕与前秦的关隘、渡口、军镇和水陆道路的可通行性继续约束使节、军队和物资的行动继续约束东晋、前燕与前秦的行动。 两条线索共同说明，政权变化之外，地方社会和空间条件仍在塑造行动者可以选择的办法。 就可说范围而言，前一判断止于《晋书》哀海西简文诸帝纪；《晋书》载记；《资治通鉴·晋纪》，本纪与编年材料主要确证政权年序和精英行动；对基层执行、疆界和人口数量的沉默不得以叙事补足；后一判断止于《水经注》；墓志、城址与十六国史辑佚，纪传中的行军路线、兵数与控制范围常被简化或夸饰，地理材料只能提供有限交叉。

沿着公元352年（年序桥接），东晋、前燕与前秦的\eventanchor{JH-DENS-352-A}{前燕、前秦与东晋围绕河北、关中和淮汉走廊轮番调兵，扩张依赖州郡供给与地方合作}指向前燕、前秦与东晋围绕河北、关中和淮汉走廊轮番调兵，扩张依赖州郡供给与地方合作，\eventanchor{JH-DENS-352-B}{东晋、前燕与前秦的朝廷和州郡围绕户籍、田租、军粮与转运的安排进入材料}则把东晋、前燕与前秦的朝廷和州郡围绕户籍、田租、军粮与转运的安排进入材料接入同一条年序。 这使读者能把下一年的军事、行政或社会变化放回尚未消散的限制中。 此处的证据支点是《晋书》哀海西简文诸帝纪；《晋书》载记；《资治通鉴·晋纪》和《水经注》；墓志、城址与十六国史辑佚；本纪与编年材料主要确证政权年序和精英行动；对基层执行、疆界和人口数量的沉默不得以叙事补足。，同时志书、诏令与本纪可显示制度设计或战争需求，不能直接换算赋额、仓储或实际负担。

在公元353年（年序桥接）的并行行动者之间，\eventanchor{JH-DENS-353-A}{前燕、前秦与东晋围绕河北、关中和淮汉走廊轮番调兵，扩张依赖州郡供给与地方合作}将前燕、前秦与东晋围绕河北、关中和淮汉走廊轮番调兵，扩张依赖州郡供给与地方合作落到时间位置，\eventanchor{JH-DENS-353-B}{东晋、前燕与前秦的流民、侨置户、军镇居民或地方家族的安置与编户问题构成社会背景}从东晋、前燕与前秦的流民、侨置户、军镇居民或地方家族的安置与编户问题构成社会背景补足资源与空间的视角。 两条线索共同说明，政权变化之外，地方社会和空间条件仍在塑造行动者可以选择的办法。 材料边界须分别保留：《晋书》哀海西简文诸帝纪；《晋书》载记；《资治通鉴·晋纪》与《水经注》；墓志、城址与十六国史辑佚所见不同，且本纪与编年材料主要确证政权年序和精英行动；对基层执行、疆界和人口数量的沉默不得以叙事补足；墓志、家族文本和侨置名目各自有选择性，不能据此统计迁徙规模或概括整体社会态度。

在公元354年（年序桥接），\eventanchor{JH-DENS-354-A}{前燕、前秦与东晋围绕河北、关中和淮汉走廊轮番调兵，扩张依赖州郡供给与地方合作}所见的前燕、前秦与东晋围绕河北、关中和淮汉走廊轮番调兵，扩张依赖州郡供给与地方合作与\eventanchor{JH-DENS-354-B}{东晋、前燕与前秦的寺院、僧侣、道士和施主网络在材料中连接都城、州郡与边地}所见的东晋、前燕与前秦的寺院、僧侣、道士和施主网络在材料中连接都城、州郡与边地并行发生在东晋、前燕与前秦的历史空间。 这使读者能把下一年的军事、行政或社会变化放回尚未消散的限制中。 证据不许把这两个观察写成全域事实：《晋书》哀海西简文诸帝纪；《晋书》载记；《资治通鉴·晋纪》受本纪与编年材料主要确证政权年序和精英行动；对基层执行、疆界和人口数量的沉默不得以叙事补足。限制，《水经注》；墓志、城址与十六国史辑佚亦受僧传、道教文本、造像记与遗址的成书或营造年代须分开，个案不能量化信仰或政策落实限制。


《晋书·载记》、帝纪与《资治通鉴》可定主线；疆域、兵数和族群归属不可平滑推定。


\subsection{355—359年：材料所见的连续压力}
穆帝初年的州郡仍在运转，征收、安置、营造和通行构成朝廷与地方互相看见的条件，也积累下一轮重组的压力。

\eventanchor{JH-DENS-355-A}{前燕、前秦与东晋围绕河北、关中和淮汉走廊轮番调兵，扩张依赖州郡供给与地方合作}和\eventanchor{JH-DENS-355-B}{东晋、前燕与前秦的关隘、渡口、军镇和水陆道路的可通行性继续约束使节、军队和物资的行动}使公元355年（年序桥接）的东晋、前燕与前秦既保留前燕、前秦与东晋围绕河北、关中和淮汉走廊轮番调兵，扩张依赖州郡供给与地方合作，也不略去东晋、前燕与前秦的关隘、渡口、军镇和水陆道路的可通行性继续约束使节、军队和物资的行动。 如此处理不是用空白填满史实，而是避免后来的转折抹去本年仍在运作的条件。 材料边界须分别保留：《晋书》哀海西简文诸帝纪；《晋书》载记；《资治通鉴·晋纪》与《水经注》；墓志、城址与十六国史辑佚所见不同，且本纪与编年材料主要确证政权年序和精英行动；对基层执行、疆界和人口数量的沉默不得以叙事补足；纪传中的行军路线、兵数与控制范围常被简化或夸饰，地理材料只能提供有限交叉。

\eventanchor{JH-DENS-356-A}{前燕、前秦与东晋围绕河北、关中和淮汉走廊轮番调兵，扩张依赖州郡供给与地方合作}在公元356年（年序桥接）提示前燕、前秦与东晋围绕河北、关中和淮汉走廊轮番调兵，扩张依赖州郡供给与地方合作仍未结案；\eventanchor{JH-DENS-356-B}{东晋、前燕与前秦的朝廷和州郡围绕户籍、田租、军粮与转运的安排进入材料}以东晋、前燕与前秦的朝廷和州郡围绕户籍、田租、军粮与转运的安排进入材料说明东晋、前燕与前秦的选择还受具体条件限制。 年序因此不把大事件当作一切的终点，而让制度、交通和社会关系继续承受压力。 此处的证据支点是《晋书》哀海西简文诸帝纪；《晋书》载记；《资治通鉴·晋纪》和《水经注》；墓志、城址与十六国史辑佚；本纪与编年材料主要确证政权年序和精英行动；对基层执行、疆界和人口数量的沉默不得以叙事补足。，同时志书、诏令与本纪可显示制度设计或战争需求，不能直接换算赋额、仓储或实际负担。

公元357年（年序桥接），东晋、前燕与前秦的\eventanchor{JH-DENS-357-A}{前燕、前秦与东晋围绕河北、关中和淮汉走廊轮番调兵，扩张依赖州郡供给与地方合作}把前燕、前秦与东晋围绕河北、关中和淮汉走廊轮番调兵，扩张依赖州郡供给与地方合作留在年序中；\eventanchor{JH-DENS-357-B}{东晋、前燕与前秦的流民、侨置户、军镇居民或地方家族的安置与编户问题构成社会背景}则从东晋、前燕与前秦的流民、侨置户、军镇居民或地方家族的安置与编户问题构成社会背景观察同一年的约束。 它们不是由年份自动推出的因果，而是让前后事件之间的资源、道路、户籍或地方网络保持可见。 证据不许把这两个观察写成全域事实：《晋书》哀海西简文诸帝纪；《晋书》载记；《资治通鉴·晋纪》受本纪与编年材料主要确证政权年序和精英行动；对基层执行、疆界和人口数量的沉默不得以叙事补足。限制，《水经注》；墓志、城址与十六国史辑佚亦受墓志、家族文本和侨置名目各自有选择性，不能据此统计迁徙规模或概括整体社会态度限制。

公元358年（年序桥接），\eventanchor{JH-DENS-358-A}{前燕、前秦与东晋围绕河北、关中和淮汉走廊轮番调兵，扩张依赖州郡供给与地方合作}和\eventanchor{JH-DENS-358-B}{东晋、前燕与前秦的寺院、僧侣、道士和施主网络在材料中连接都城、州郡与边地}共同避免把东晋、前燕与前秦压缩为单一都城事件：前者关注前燕、前秦与东晋围绕河北、关中和淮汉走廊轮番调兵，扩张依赖州郡供给与地方合作，后者关注东晋、前燕与前秦的寺院、僧侣、道士和施主网络在材料中连接都城、州郡与边地。 这里保留的是可核的连续性，而非对基层反应、疆界或人数的无据补写。 就可说范围而言，前一判断止于《晋书》哀海西简文诸帝纪；《晋书》载记；《资治通鉴·晋纪》，本纪与编年材料主要确证政权年序和精英行动；对基层执行、疆界和人口数量的沉默不得以叙事补足；后一判断止于《水经注》；墓志、城址与十六国史辑佚，僧传、道教文本、造像记与遗址的成书或营造年代须分开，个案不能量化信仰或政策落实。

把公元359年（年序桥接）放回连续叙事，\eventanchor{JH-DENS-359-A}{前燕、前秦与东晋围绕河北、关中和淮汉走廊轮番调兵，扩张依赖州郡供给与地方合作}保留前燕、前秦与东晋围绕河北、关中和淮汉走廊轮番调兵，扩张依赖州郡供给与地方合作，\eventanchor{JH-DENS-359-B}{东晋、前燕与前秦的关隘、渡口、军镇和水陆道路的可通行性继续约束使节、军队和物资的行动}让东晋、前燕与前秦的关隘、渡口、军镇和水陆道路的可通行性继续约束使节、军队和物资的行动继续约束东晋、前燕与前秦的行动。 如此处理不是用空白填满史实，而是避免后来的转折抹去本年仍在运作的条件。 此处的证据支点是《晋书》哀海西简文诸帝纪；《晋书》载记；《资治通鉴·晋纪》和《水经注》；墓志、城址与十六国史辑佚；本纪与编年材料主要确证政权年序和精英行动；对基层执行、疆界和人口数量的沉默不得以叙事补足。，同时纪传中的行军路线、兵数与控制范围常被简化或夸饰，地理材料只能提供有限交叉。


《晋书》、志书与《资治通鉴》只显示制度意图，不能证明每一州郡照章施行。
